\section{Related Work}
\label{sec:related}

\para{Attention in natural language processing.}
The attention mechanism was introduced by \citet{bahdanau2014neural} for neural machine translation and elevated to architectural primacy by \citet{vaswani2017attention}, who showed that self-attention alone---without recurrence or convolution---achieves state-of-the-art translation quality. \citet{devlin2019bert} extended this insight to bidirectional pretraining, producing universal language representations that transfer across tasks. These advances established attention as the dominant paradigm in NLP, but they also created a perception that attention is inherently a language-processing technique.

\para{Attention in computer vision.}
\citet{dosovitskiy2021image} challenged this perception by applying a pure transformer to image classification, treating image patches as tokens. The Vision Transformer (\vit) matches or exceeds convolutional networks when pretrained on large datasets, though it underperforms on small data due to the absence of built-in spatial inductive biases. \citet{guo2022attention} survey the subsequent proliferation of attention mechanisms in vision, covering channel attention, spatial attention, and their combinations across classification, detection, segmentation, and generation tasks. More recently, \citet{peebles2023scalable} replaced the U-Net backbone in diffusion models with transformers, and this architecture now underlies leading image and video generation systems~\citep{esser2024stable}.

\para{Attention in other modalities.}
The attention mechanism has spread well beyond text and images. \citet{radford2022whisper} use transformer-based encoder-decoder attention for robust speech recognition across languages. \citet{huang2018music} apply self-attention with relative positional encodings to music generation. \citet{velickovic2018graph} introduce graph attention networks that learn to weight neighbor contributions in graph-structured data, directly applicable to social and citation networks. \citet{brauwers2022survey} provide a comprehensive cross-domain taxonomy, showing that the same attention formulation---feature extraction, query computation, relevance scoring, and weighted aggregation---applies uniformly across NLP, vision, audio, recommendation, and graph domains.

\para{Attention in recommendation and ranking.}
Self-attention has also reshaped recommendation systems. \citet{kang2018self} propose \textsc{SASRec}, which applies self-attention to sequential recommendation and outperforms RNN and CNN baselines. \citet{linkedin2025ranking} describe LinkedIn's production ranking system, which uses a modified transformer with only seven features to achieve state-of-the-art content ranking at scale. These systems allocate user attention---deciding which content appears in a feed---using the same mathematical mechanism that allocates computational attention within a neural network.

\para{The attention economy.}
\citet{simon1971designing} articulated the foundational insight: ``a wealth of information creates a poverty of attention and a need to allocate that attention efficiently among the overabundance of information sources.'' This economic principle directly parallels the computational attention mechanism: both involve allocating limited capacity to the most relevant sources. Recent work on attention economics identifies attention as a ``universal symbolic currency'' on social media and beyond~\citep{davenport2001attention}. Our work bridges these two senses of attention by showing they share not just a metaphorical but a mathematical connection.

\para{Positioning our work.}
Unlike prior surveys that document attention's spread across domains~\citep{brauwers2022survey,guo2022attention}, we conduct controlled experiments comparing attention-based and non-attention models across modalities using identical architectures. Unlike work on specific applications~\citep{kang2018self,velickovic2018graph}, we explicitly map the mathematical structure shared by transformer attention, PageRank, and collaborative filtering. Our goal is not to propose a new attention variant but to establish attention as a universal computational principle with formal connections to internet infrastructure.
